\subsection*{Warmth}
\paragraph{Sentences that raise the score.}
\begin{itemize}
  \item ``They are taught how to show respect for different opinions and listen to the views of others.''
  \item ``They hold leaders strictly to account for the quality of the school's provision and are highly effective.''
  \item ``They listen to others and express their views in a well-managed environment.''
  \item ``Sixth-form students are proud role models and make a highly positive contribution to the school and local community.''
  \item ``Those responsible for governance know the school's context and provide support and professional challenge where necessary to ensure that the school's high standards are maintained.''
  \item ``In lessons, pupils are diligent, focused, and eager to learn.''
  \item ``They listen to staff and take into account the pressures on them.''
  \item ``Pupils are encouraged to contribute to the school community.''
  \item ``In lessons, pupils listen respectfully to teachers and each other.''
  \item ``The school's motto is `excellence with care.' Pupils testify to how this works in practice.''
  \item ``In the secondary phase, pupils receive regular careers information and guidance.''
  \item ``They move calmly between classes and wear their uniforms smartly, meeting the school's high standards.''
\end{itemize}
\paragraph{Sentences that lower the score.}
\begin{itemize}
  \item ``However, some pupils, particularly some disadvantaged pupils, are absent too often.''
  \item ``Many pupils feel confident they have a `trusted adult' they can go to.''
  \item ``(Information for the school and appropriate authority)  In some subjects, the work given to some pupils does not enable them to develop their ideas and extend their answers.''
  \item ``Pupils take part in a range of extra-curricular activities to develop their talents and interests, for example basketball, choir, photography, and cookery.''
  \item `` Sometimes, the work given to pupils, including those with SEND, does not enable them to achieve the aims and ambition of the curriculum.''
  \item ``However, in some subjects, the work given to some pupils does not enable them to develop their ideas and extend their answers as well as they could.''
  \item ``Missing out on important learning means that these pupils do not achieve as well as they might.''
  \item ``Pastoral staff work closely with these pupils and their families.''
  \item ``They check closely to ensure that disadvantaged pupils have the support they need.''
  \item ``The work given to pupils does not always give them opportunities to develop deeper thinking, such as reasoning and justifying.''
  \item ``Leaders monitor closely pupils' participation in the school's extra-curricular offer.''
  \item ``The school recognises the importance of ensuring that pupils, particularly disadvantaged pupils, and pupils with SEND, gain wider opportunities beyond the curriculum.''
\end{itemize}
\paragraph{The strongest longer phrases.} Raising: ``high expectations behaviour''; ``early career teachers''; ``provide effective support''; ``pupils attend school regularly''; ``school high expectations''; ``positive attitudes learning''; ``meet pupils needs''; ``pupils need help''. Lowering: ``work given pupils''; ``support pupils learning''; ``develop leadership skills''; ``expectations pupils achieve''; ``inspection headteacher school''; ``pupils kept safe''; ``broad ambitious curriculum''; ``ambitious curriculum pupils''.
\subsection*{Teaching}
\paragraph{Sentences that raise the score.}
\begin{itemize}
  \item ``Those responsible for governance know the school's context and provide support and professional challenge where necessary to ensure that the school's high standards are maintained.''
  \item ``Leaders are supported and challenged by dedicated and knowledgeable governors, who share their vision of high standards in a supportive school.''
  \item ``In a few subjects, however, pupils are not remembering knowledge deeply in the long term.''
  \item ``The school's motto is `excellence with care.' Pupils testify to how this works in practice.''
  \item ``They lead charity initiatives and community-based projects to instil a keen sense of responsibility and empathy for others.''
  \item ``In lessons pupils settle to their work quickly and show high levels of respect to their teachers and to other pupils.''
  \item ``They are taught how to show respect for different opinions and listen to the views of others.''
  \item ``In lessons, most pupils focus on their learning and learn well.''
  \item ``They carefully choose the core knowledge they want pupils to know and remember in the long term.''
  \item ``Staff, pupils, and parents share a strong sense of pride in belonging to this school's community.''
  \item ``Leaders have ensured that all aspects of the national curriculum are taught in Years to Work with primary schools helps new pupils to settle in quickly.''
  \item ``Teachers help pupils to embed learning in their long-term memory.''
\end{itemize}
\paragraph{Sentences that lower the score.}
\begin{itemize}
  \item ``The personal, social and health education (PSHE) curriculum is designed to support pupils' wider development.''
  \item ``Pupils take part in a range of extra-curricular activities to develop their talents and interests, for example basketball, choir, photography, and cookery.''
  \item ``On other occasions, adults do not support pupils to complete work to a high standard.''
  \item ``Judgements in this report are based on the current inspection framework and also reflect changes that have happened at any point since the last graded inspection.''
  \item ``The school offers pupils a range of ways to develop their wider skills and talents.''
  \item ``The school has carefully designed the personal, social and health education curriculum so that it supports pupils' wider development.''
  \item ``It does so, not only by having high expectations for pupils, but by supporting their individual aspirations and talents.''
  \item ``The curriculum in personal, social and health education is designed to help pupils understand important ideas.''
  \item ``Well-planned personal, social, health education ensures that pupils learn about relevant and sensitive issues, including personal safety, relationships, sexual health and consent.''
  \item ``However, some pupils, particularly some disadvantaged pupils, are absent too often.''
  \item ``They have an accurate understanding of what remains to be done to achieve this for all pupils.''
  \item ``Pastoral staff work closely with these pupils and their families.''
\end{itemize}
\paragraph{The strongest longer phrases.} Raising: ``early career teachers''; ``school high expectations''; ``high expectations behaviour''; ``important knowledge pupils''; ``behave lessons school''; ``pupils need additional''; ``positive attitudes learning''; ``key knowledge skills''. Lowering: ``support pupils learning''; ``develop leadership skills''; ``school community pupils''; ``pupils kept safe''; ``high expectations pupils achieve''; ``ambitious curriculum pupils''; ``expectations pupils achieve''; ``school pupils enjoy''.
